\section{A direct position integral for the creation block}
\label{sec:position}
We now remove the auxiliary spacetime and group cutoffs from a specific
position pairing. The proper-time Gaussian and the actual observer
threshold remain fixed. The output is tested by an original seed, not
by a second invariant wavefunction integrated over the full gauge orbit.
This distinction avoids the group-volume divergence noted after
(5.8) of~\cite{KL}. Transposition to test vectors is standard; the
domain estimates below supply the additional input needed here.
In particular, the invariant-domain requirements in
refined algebraic quantization~\cite[Section II]{GM} cannot be replaced
by formal commutation with the constraints. The distinction between
operators on invariants and co-invariants is also discussed
in~\cite[Section 5.1]{HM}.

\subsection{Test vectors and observer time}
Let $\mathcal P_N$ be the finite span of actual product seeds
\begin{equation}
 \phi(y)=u_{0,f_L}(x_L)\,u_{n_R,f_R}(x_R)\,v_N,
 \qquad f_L,f_R\in\Cc((s_*,\infty)),
 \label{eq:position-seeds}
\end{equation}
where $v_N$ is a finite linear combination of normalized matter occupation
vectors and $n_R+J_{\mathrm m}=0$ term by term. These are the small
total-$K=0$ seeds used in Section~\ref{sec:physical}; $B\mathcal P_N$
is dense in the physical sector. The case $N=0$ has $n_R=0$.

Write $E_R(u)=e^{iuH_R}$ on the kinematical observer space. It induces
a strongly continuous physical unitary $\mathsf U_N(u)$:
in the identified spectrum it multiplies by $e^{iu m(a)}$, where
$a$ is the actual right-observer label, with $a=d$ for $N=0$.
Thus, on these seeds,
\begin{equation}
 \mathsf U_N(u)B\psi=B E_R(u)\psi,\qquad
 \norm{\mathsf U_N(u)k}_{\FP}=\norm k_{\FP}.
 \label{eq:physical-observer-time}
\end{equation}
This construction uses the already identified state space and does not
assume a local field. Since the mass phase preserves compact spectral
support, $E_R(u)\mathcal P_N\subset\mathcal P_N$ for every finite $u$.

For a finite-mode matter vector define the ordinary, finite sum
\begin{equation}
 \mathfrak a_\nu(x)v=\sum_n\psi^\nu_n(x)a_n v.
 \label{eq:finite-annihilation}
\end{equation}
The sum is finite on this domain. Fixed-$n$ mode asymptotics give
$\sup_x|\psi_n^\nu(x)|<\infty$; more precisely
$|\psi_n^\nu(t,\theta)|\le c_{n,\nu}(\cosh t)^{-1/2}$.
Consequently
\begin{equation}
 \norm{\mathfrak a_\nu(x)v}\le
 \sum_{n\in F_v}c_{n,\nu}\norm{a_nv}=:c_v.
 \label{eq:finite-annihilation-bound}
\end{equation}
No uniform estimate as the occupied mode set $F_v$ grows is asserted.

For $\phi\in\mathcal P_{N+1}$ set
\begin{equation}
 \zeta_{\phi,u}(y)=
 \mathfrak a_\nu(x_R)\,[E_R(u)\phi](y)\in\mathcal K_N.
 \label{eq:transposed-test}
\end{equation}
Although $\zeta_{\phi,u}$ need not obey the observer positive-frequency
condition, it is an ambient orbit test. For a single product in
\eqref{eq:position-seeds}, its pointwise norm is at most
$c_v|u_{0,f_L}(x_L)u_{n_R,e^{ium}f_R}(x_R)|$.
The norm inequality survives untwisting because the matter representation
is unitary.

\begin{lemma}[Polynomial orbit bound for the transposed test]
For every fixed $\phi\in\mathcal P_{N+1}$ there is a constant $C_\phi$
such that
\begin{equation}
 \zeta_{\phi,u}\in\cD_{\mathrm{orb}},\qquad
 \norm{a_{\zeta_{\phi,u}}}_2\le C_\phi(1+|u|)^2.
 \label{eq:transposed-orbit-bound}
\end{equation}
For every fixed seed $\psi\in\mathcal P_N$,
$\norm{a_{E_R(u)\psi}}_2\le C_\psi(1+|u|)^2$ as well.
\end{lemma}
\begin{proof}
Fix $\alpha=1/2$ in \eqref{eq:product-orbit-bound}. The two-ended packet
expansion used in Section~\ref{sec:spectral}, with two integrations by
parts in its compact mass variable, yields
\[
 \sup_t(1+|t|)^2
 \left|\sqrt{2\pi\cosh t}\,u_{n,e^{ium}f}(t,0)\right|
 \le C_{n,\supp f}\sum_{j=0}^2\norm{\partial_s^j(e^{ium(s)}f(s))}_1
 \le C_f(1+|u|)^2.
\]
The same estimate on the bounded-$t$ part follows from smoothness of
the modes. The exponent $2$ is enough because
$\int(1+|t|)^{-3/2}\dd t<\infty$; the radial $L^2$ norm is unchanged
by the mass phase. Apply \eqref{eq:product-orbit-bound} and the
pointwise domination following \eqref{eq:transposed-test}.
Finite sums satisfy the triangle inequality for the orbit seminorm.
\end{proof}

\subsection{The uncut integral and its identification}
Let
\begin{equation}
 g_{\sigma,\tau}(u)=
 \frac{1}{\sqrt{2\pi}\sigma}e^{-(u-\tau)^2/(2\sigma^2)}.
\end{equation}
For an arbitrary physical input $k\in\cH_N^{\max,w}(s_*)$ and
$\phi\in\mathcal P_{N+1}$, define directly
\begin{equation}
 \mathcal I_N^+(\phi,k)=
 \int_{\mathbb R}g_{\sigma,\tau}(u)\dd u
 \int_Y
 \ip{\zeta_{\phi,u}(y)}
    {[R\mathsf U_N(u)k](y)}\dd y.
 \label{eq:direct-position-creation}
\end{equation}
Inner products are conjugate-linear in the first entry. Both observer
legs of the test and the physical input already obey the threshold.
The two factors $E_R(u)$ and $\mathsf U_N(u)$ therefore produce
$e^{-iu[m(a)-m(b)]}$ in a finite-mode matrix element, with precisely
the sign in \eqref{eq:gaussian-time}.

\begin{theorem}[Direct creation pairing]
Integral \eqref{eq:direct-position-creation} is absolutely convergent
after taking the norm of each pointwise Hilbert pairing. It satisfies
\begin{equation}
 \mathcal I_N^+(\phi,k)
 =\ip{B\phi}{\mathsf C_N k}_{\FP},
 \qquad
 |\mathcal I_N^+(\phi,k)|
 \le\sqrt{(N+1)\mathfrak M}\,
       \norm{B\phi}_{\FP}\norm k_{\FP}.
 \label{eq:direct-position-identified}
\end{equation}
Here $\mathsf C_N$ is the previously constructed Gaussian creation
block at the same $\sigma,\tau,s_*$. Thus the direct formula determines
the same physical output for every physical input.
\label{thm:direct-position}
\end{theorem}
\begin{proof}
First use \eqref{eq:position-pairing} and
\eqref{eq:transposed-orbit-bound}, before invoking any field:
\begin{align}
 &\int_{\mathbb R}g_{\sigma,\tau}(u)\dd u
 \int_Y\norm{\zeta_{\phi,u}(y)}
          \norm{[R\mathsf U_N(u)k](y)}\dd y\notag\\
 &\hspace{12mm}\le C_\phi\norm k_{\FP}
       \int_{\mathbb R}g_{\sigma,\tau}(u)(1+|u|)^2\dd u<\infty.
 \label{eq:direct-absolute-bound}
\end{align}
Jointly measurable representatives can be chosen: approximate $k$
by physical seed images, use strong continuity of $\mathsf U_N$,
and use the continuity of $R$ into local $L^1$ on each compact
time interval and compact subset of $Y$. The resulting pairing is
independent of representatives. This also proves continuity in $k$.

It remains to identify this independently defined integral.
Take first $k=B\psi$, $\psi\in\mathcal P_N$, and put
$\psi_u=E_R(u)\psi$. The stronger, pre-cancellation estimate is
\begin{align}
 &\int_G\dd g\int_Y\dd y\,
     \norm{\zeta_{\phi,u}(y)}\norm{[U_N(g)\psi_u](y)}
   =\int_Z a_{\zeta_{\phi,u}}a_{\psi_u}\dd\mu\notag\\
 &\hspace{12mm}\le C_\phi C_\psi(1+|u|)^4 .
 \label{eq:triple-fubini-bound}
\end{align}
The equality follows by the same nonnegative coarea and Haar changes of
variables as \eqref{eq:fp-ga}. Its Gaussian integral is finite, so
the $u$, $g$, and $y$ integrals may now be interchanged.
In particular the group-first integral is
$R\mathsf U_N(u)B\psi$, not a finite-cutoff non-invariant state.

For fixed $g$, the $u,y$ integral equals
$\ip{\phi}{(1_L\otimes C_N)U_N(g)\psi}_{\rm kin}$.
Here is a finite-mode justification that does not assume an infinite
mode interchange. The test-side annihilation in
\eqref{eq:finite-annihilation} leaves finitely many matter vectors.
Only their overlaps with the transformed input matter vector contribute.
Integration in $\theta_L$ selects the test's left mode, and integration
in $\theta_R$ selects finitely many input observer modes by
$n_b=n_a+n$. At fixed $g$ their mass coefficients are the original
compact packets times smooth principal-series matrix coefficients.
Each surviving radial matrix element is therefore the absolutely
convergent triple-mode integral \eqref{eq:local-mode}.
The Gaussian integral supplies \eqref{eq:gaussian-time}.
The same finite selection follows for the bounded $C_N$ by
approximating the transformed input in its original-mode Hilbert basis;
the forbidden components have identically zero matrix elements.
Thus no uniform large-angular-label bound has been assumed.

Integrating over $g$ now gives the previously defined form
\eqref{eq:weak-local-definition}, hence
$\ip{B\phi}{\mathsf C_NB\psi}_{\FP}$.
Finally $B\mathcal P_N$ is dense, \eqref{eq:direct-absolute-bound}
is continuous in its physical input, and $\mathsf C_N$ is bounded.
Approximation extends the identity to all $k$ and proves the physical
norm bound in \eqref{eq:direct-position-identified}.
\end{proof}

The identification also proves that \eqref{eq:direct-position-creation}
vanishes for a physically null test seed. Representative independence
and the bound by $\norm{B\phi}$ follow from the identity; they do not
follow from a bound involving $C_\phi$ alone.

There is a useful cutoff consequence. Let $0\le\xi_j\le1$ be compact
position cutoffs tending pointwise to one. For each fixed test $\phi$,
the error from inserting $\xi_j(y)$ in
\eqref{eq:direct-position-creation} is at most
\begin{equation}
 \norm k_{\FP}
 \int_{\mathbb R}g_{\sigma,\tau}(u)
       \norm{a_{(1-\xi_j)\zeta_{\phi,u}}}_2\dd u
 \longrightarrow0.
 \label{eq:position-cutoff-error}
\end{equation}
At fixed $u$ this is dominated convergence in the orbit norm; the
polynomial bound supplies the Gaussian majorant. A proper-time cutoff
has the analogous tail bound from \eqref{eq:direct-absolute-bound}.
Hence removal of these cutoffs converges uniformly on the physical
input unit ball for a fixed test, with no exchange of conditionally
convergent integrals.

\subsection{What is still not a direct position theorem}
Theorem~\ref{thm:direct-position} identifies the projected, time-smeared
creation pairing on an explicit dense physical test family and arbitrary
physical inputs. It does not assert that the raw product
$\varphi^+(x_R)Rk(y)$ is a pointwise $\mathcal K_{N+1}$ vector,
or that raw multiplication preserves observer positive frequency.
The test-side finite sum is the essential domain choice.

The annihilation block is already the physical adjoint
$\mathsf C_N^*$, but its unrestricted direct position formula requires
an additional argument. Transposing it would put
$\mathfrak a_\nu^\dagger(x_R)$ on the test: even for the matter vacuum
this is an infinite-mode point-field expression, so
\eqref{eq:finite-annihilation-bound} is unavailable.
Nor does the present proof remove the Gaussian, establish the
$s_*=0$ field extension, or justify integrating two physical invariant
states over the full unreduced position space.
